\scanpage{31}
\mspage{3}

est cocartésien -- ou encore, si $Y,X,X^*$ sont connexes, et qu'on choisit un ``pt base'' de $\Pi_1\unsure{\mathcal V_{Y,X}^*|V_{Y,X}^*}$ (i.e. par définition, un rev. universel de $\unsure{\mathcal V_{Y,X}^*|V_{Y,X}^*}$) (d'où des pts bases de $\Pi_1\unsure{\mathcal V_{Y,X}|V_{Y,X}}$, \deleted{\unsure{et un point} de $\Pi_1X^*)$}\inserted{$\Pi_1X^*,\Pi_1X)$}, qu'on a un isom. \inserted{\unsure{can.}} de groupes fondamentaux
\[
(9)\qquad \unsure{\pi|\Pi}_1(X)\simeq\unsure{\pi|\Pi}_1(X^*)\wedge_{\unsure{\pi|\Pi}_1\unsure{\mathcal V_{Y,X}^*|V_{Y,X}^*}}\underbrace{\unsure{\pi|\Pi}_1(\unsure{\mathcal V_{Y,X}|V_{Y,X}})}_{\simeq\unsure{\pi|\Pi}_1(Y)},
\]
où $\unsure{\pi|\Pi}_1(\unsure{\mathcal V_{Y,X}|V_{Y,X}})$ est isom. \inserted{\unsure{extérieurem\textsuperscript{t}}} à $\unsure{\pi|\Pi}_1(Y)$.

\deleted{\unsure{relatif} à un pt base choisi de}

Pour expliciter $\unsure{\pi|\Pi}_1(X)$ en termes des données ``élémentaires'', \deleted{il reste} dont $\unsure{\pi|\Pi}_1(Y)$ \inserted{et} $\unsure{\pi|\Pi}_1(X^*)$, \deleted{et} il \unsure{resterait} encore à expliciter la structure de $\unsure{\pi|\Pi}_1(\unsure{\mathcal V_{Y,X}^*|V_{Y,X}^*})$, qui s'envoie dans l'un et dans l'autre, donnant lieu justement au diagramme

\margin{\unsure{cette flèche définie moy\textsuperscript{t}} d'un choix suppl. (pt base de $Y$ et classe de chemins de $x$ vers $y$)}

(10)
\diagram{\begin{tikzcd}[ampersand replacement=\&]
\unsure{\pi|\Pi}_1(\unsure{\mathcal V_{Y,X}^*|V_{Y,X}^*})\arrow[r]\arrow[d]\&\unsure{\pi|\Pi}_1(\deleted{\unsure{X}}Y)\\
\unsure{\pi|\Pi}_1(X^*)\&
\end{tikzcd}}

qui exprime (8). C'est ici que l'hypothèse

\margin{30}
